\documentclass[12pt]{article}
\usepackage{fullpage}
\usepackage[T1]{fontenc}
\usepackage[utf8]{inputenc}
\usepackage{amsmath,amssymb,amsthm}
\usepackage{hyperref}

\newcommand{\GL}{\mathrm{GL}}
\newcommand{\oo}{\mathfrak{o}}
\newcommand{\pp}{\mathfrak{p}}
\newcommand{\qq}{\mathfrak{q}}
\newcommand{\Wmod}{\mathcal{W}}
\newcommand{\diag}{\operatorname{diag}}
\DeclareMathOperator{\Cc}{C_c^\infty}

\newtheorem{theorem}{Theorem}
\newtheorem{lemma}{Lemma}
\newtheorem{proposition}{Proposition}

\title{Problem 2: A uniform test vector for the local Rankin--Selberg integral}
\author{}
\date{}

\begin{document}
\maketitle

\begin{theorem}
The answer is \emph{yes}. Let $\Pi$ be a generic irreducible admissible
representation of $\GL_{n+1}(F)$, realized in $\Wmod(\Pi,\psi^{-1})$. Then there
exists $W\in\Wmod(\Pi,\psi^{-1})$, depending only on $\Pi$, such that for every
generic irreducible admissible representation $\pi$ of $\GL_n(F)$ there is a vector
$V\in\Wmod(\pi,\psi)$ for which the local Rankin--Selberg integral
\[
  \Psi(s,W,V):=\int_{N_n\backslash \GL_n(F)} W(\diag(g,1)\,u_Q)\,V(g)\,|\det g|^{s-\frac12}\,dg
\]
is finite and nonzero for all $s\in\mathbb C$. (Here, as in the statement,
$\qq$ is the conductor ideal of $\pi$, $Q$ generates $\qq^{-1}$, and
$u_Q=I_{n+1}+Q\,E_{n,n+1}$.)
\end{theorem}

Throughout, $F$ is a non-archimedean local field with ring of integers $\oo$,
maximal ideal $\pp$, and residue field of size $q_F$; $\psi$ has conductor $\oo$,
i.e.\ $\psi|_{\oo}=1$ and $\psi|_{\pp^{-1}}\neq1$. We write $K=\GL_n(\oo)$ for the
maximal compact subgroup, $N_r\subset\GL_r$ for the upper triangular unipotent
subgroup, $P_r\subset\GL_r$ for the mirabolic subgroup (matrices whose bottom row
is $(0,\dots,0,1)$), and $e_r=(0,\dots,0,1)^{\mathsf T}$. A function on a
homogeneous space is ``$(N_r,\chi)$-equivariant'' if $f(ng)=\chi(n)f(g)$ for
$n\in N_r$, and $\Cc(N_r\backslash G,\chi)$ denotes the smooth such functions with
support compact modulo $N_r$.

\paragraph{Overview and the key idea.}
The single delicate point in a problem of this shape is that the vector $W$ must be
fixed \emph{before} $\pi$ is given, while $\pi$ (hence its conductor exponent $c$
and the generator $Q$ with $v(Q)=-c$) varies without bound. One is therefore not
free to choose $W$ adapted to the conductor of $\pi$.

We resolve this by choosing $W$ so that the restriction
$W|_{P_{n+1}}$ is a fixed function \emph{supported on the image of the mirabolic
$P_n$}, more precisely on $N_{n+1}\,(P_n\cap K)$. The point is that for
$g$ in the mirabolic $P_n$ the $(n,n)$-entry is $g_{nn}=1$; consequently the only
effect of the shift $u_Q$ on the integrand is to multiply it by the nonzero
\emph{constant} $\psi(-Q)$. After this collapse the integral involves only the
restriction of $V$ to $P_n\cap K$, which can be prescribed freely by the Kirillov
model of $\pi$. No information about the conductor of $\pi$ is needed to construct
$W$, and the resulting integral is constant in $s$, hence finite and nonzero for
all $s$.

The proof has three steps.
\begin{itemize}
\item[(A)] We record the effect of $\diag(g,1)u_Q$ in mirabolic coordinates and the
identification of $\Cc(N_{n+1}\backslash P_{n+1},\psi^{-1})$ with
$\Cc(N_n\backslash\GL_n,\psi^{-1})$.
\item[(B)] We choose, once and for all, a single vector $W$ depending only on
$\Pi$, with $W|_{P_{n+1}}$ supported on $N_{n+1}(P_n\cap K)$, and we compute
$\Psi(s,W,V)$ explicitly: it collapses to $\psi(-Q)\cdot M(V)$ with
$M(V):=\int_{(N_n\cap K)\backslash(P_n\cap K)}V(p)\,dp$, independent of $s$.
\item[(C)] We produce $V\in\Wmod(\pi,\psi)$ with $M(V)\neq0$ using only the
surjectivity of the Kirillov model onto $\Cc(N_n\backslash P_n,\psi)$.
\end{itemize}

\section{Mirabolic coordinates and the effect of $u_Q$}
\label{sec:mirabolic}

Embed $\GL_n\hookrightarrow\GL_{n+1}$ by $g\mapsto\diag(g,1)$. Both $\diag(g,1)$
and $u_Q=I_{n+1}+Q\,E_{n,n+1}$ lie in $P_{n+1}$: the bottom row of $u_Q$ is
$e_{n+1}^{\mathsf T}=(0,\dots,0,1)$. Hence $\diag(g,1)u_Q\in P_{n+1}$ for all
$g\in\GL_n$.

Identify $P_{n+1}=\GL_n\ltimes F^n$ via
\[
  (h,v)\longleftrightarrow\begin{pmatrix} h & v\\ 0 & 1\end{pmatrix},
  \qquad h\in\GL_n,\ v\in F^n\ (\text{column}),
\]
with multiplication $(h,v)(h',v')=(hh',hv'+v)$. In these coordinates
$\diag(g,1)=(g,0)$, and since right multiplication by $u_Q$ adds $Q$ times the
$n$-th column of $\diag(g,1)$ (namely $Q\,g\,e_n$) into the last column,
\begin{equation}\label{eq:uQ}
  \diag(g,1)\,u_Q=(g,\;Q\,g e_n),\qquad g e_n=\text{$n$-th column of }g.
\end{equation}

The subgroup $N_{n+1}\subset P_{n+1}$ consists of the pairs $(u,y)$ with $u\in N_n$
and $y\in F^n$ arbitrary; its generic character $\psi$ evaluates on $(u,y)$ as
$\psi$ of the standard sum of superdiagonal entries of $u$ times $\psi(y_n)$
(the $(n,n+1)$-entry of $\diag(u,1)\cdot(I,y)$ is $y_n$). Thus
$\psi^{-1}((I,y))=\psi(-y_n)$.

\begin{lemma}\label{lem:Cc-iso}
The restriction map $f\mapsto F$, $F(g):=f((g,0))$, is a linear isomorphism
\[
  \Cc(N_{n+1}\backslash P_{n+1},\psi^{-1})\;\xrightarrow{\ \sim\ }\;
  \Cc(N_n\backslash \GL_n,\psi^{-1}),
\]
and for every $f$ in the source and every $g\in\GL_n$,
\begin{equation}\label{eq:twist}
  f\bigl(\diag(g,1)\,u_Q\bigr)=\psi(-Q\,g_{nn})\,F(g),
\end{equation}
where $g_{nn}$ is the $(n,n)$-entry of $g$.
\end{lemma}

\begin{proof}
For $(h,v)\in P_{n+1}$ we have $(h,v)=(I,v)(h,0)$ with $(I,v)\in N_{n+1}$; left
$(N_{n+1},\psi^{-1})$-equivariance of $f$ gives
$f((h,v))=\psi^{-1}((I,v))\,f((h,0))=\psi(-v_n)\,F(h)$,
so $f$ is determined by $F$. The function $F$ is left
$(N_n,\psi^{-1})$-equivariant: for $u\in N_n$, $(u,0)\in N_{n+1}$ and
$F(ug)=f((u,0)(g,0))=\psi^{-1}((u,0))F(g)=\psi^{-1}(u)F(g)$, where
$\psi^{-1}((u,0))$ is the standard character of $u\in N_n$. Smoothness and the
compact-support condition correspond under the identification
$N_{n+1}\backslash P_{n+1}=N_n\backslash\GL_n$, $(h,0)\mapsto h$. Conversely any
$F\in\Cc(N_n\backslash\GL_n,\psi^{-1})$ defines $f((h,v)):=\psi(-v_n)F(h)$, which is
well defined and lies in the source: the consistency
$f((u,0)(h,v))=\psi^{-1}(u)f((h,v))$ holds because $(u,0)(h,v)=(uh,uv)$ and
$\psi^{-1}((u,0))=\psi^{-1}(u)$ while $(uv)_n=v_n$ (the last row of $u\in N_n$ is
$e_n^{\mathsf T}$, so $uv$ has the same last coordinate as $v$). This proves the
isomorphism.

For \eqref{eq:twist}, use \eqref{eq:uQ}: $\diag(g,1)u_Q=(g,Q\,ge_n)$, and the last
coordinate of $Q\,g e_n$ is $Q\,(g e_n)_n=Q\,g_{nn}$. Hence
$f((g,Q\,ge_n))=\psi(-Q\,g_{nn})F(g)$.
\end{proof}

\section{The fixed vector $W$ and the collapse of the integral}
\label{sec:chosenW}

\subsection{The image of the restriction $W\mapsto W|_{P_{n+1}}$}

We use the following standard structural fact about generic representations.

\begin{lemma}[Mirabolic model of $\Pi$]\label{lem:miraPi}
Let $\Pi$ be a generic irreducible admissible representation of $\GL_{n+1}(F)$ in
$\Wmod(\Pi,\psi^{-1})$. The restriction map $W\mapsto W|_{P_{n+1}}$ is injective,
and its image contains $\Cc(N_{n+1}\backslash P_{n+1},\psi^{-1})$. In particular,
for every $f_0\in\Cc(N_{n+1}\backslash P_{n+1},\psi^{-1})$ there exists
$W\in\Wmod(\Pi,\psi^{-1})$ with $W|_{P_{n+1}}=f_0$.
\end{lemma}

\begin{proof}
This is the Bernstein--Zelevinsky theory of derivatives, equivalently the Kirillov
model of $\Pi$ restricted to the mirabolic. Concretely, the restriction of $\Pi$ to
$P_{n+1}$ has a filtration whose top quotient is the ``highest derivative'' piece
$(\Pi)^{(n+1)}$, and the corresponding submodule of $\Wmod(\Pi,\psi^{-1})|_{P_{n+1}}$
is the compactly induced representation
$\operatorname{ind}_{N_{n+1}}^{P_{n+1}}\psi^{-1}=\Cc(N_{n+1}\backslash P_{n+1},\psi^{-1})$;
for generic $\Pi$ this top piece is nonzero, and the restriction map to $P_{n+1}$ is
injective because $\Wmod(\Pi,\psi^{-1})$ is generated under $\GL_{n+1}(F)$ by its
$P_{n+1}$-restriction and $\Pi$ is irreducible. See
Bernstein--Zelevinsky, \emph{Ann.\ Sci.\ ENS} 10 (1977), \S5 (esp.\ Thm.\ 5.20 and
the description of the highest derivative), and
Jacquet--Piatetski-Shapiro--Shalika, \emph{Amer.\ J.\ Math.} 105 (1983), \S2.
\end{proof}

\subsection{Definition of $W$}

Let $P_n\cap K=\{p\in P_n:\ p\in M_n(\oo),\ \det p\in\oo^\times\}$, the integral
points of the mirabolic; it is a compact open subgroup of $P_n$, and its bottom row
is $e_n^{\mathsf T}$. Note that $N_n\cap K=N_n(\oo)\subseteq P_n\cap K$ and the
quotient $(N_n\cap K)\backslash(P_n\cap K)$ is compact.

Define
\begin{equation}\label{eq:F0}
  F_0\in\Cc(N_n\backslash\GL_n,\psi^{-1}),\qquad
  F_0(u\,p):=\psi^{-1}(u)\quad(u\in N_n,\ p\in P_n\cap K),\qquad
  F_0\equiv0\ \text{off}\ N_n(P_n\cap K).
\end{equation}

\begin{lemma}\label{lem:F0wd}
$F_0$ is a well-defined element of $\Cc(N_n\backslash\GL_n,\psi^{-1})$, depending
only on $n$ and $F$ (not on $\pi$ or $Q$).
\end{lemma}

\begin{proof}
\emph{Well-defined.} Suppose $u\,p=u'\,p'$ with $u,u'\in N_n$ and $p,p'\in P_n\cap K$.
Then $u'^{-1}u=p'p^{-1}\in N_n\cap(P_n\cap K)=N_n\cap K=N_n(\oo)$. Since the
superdiagonal entries of any element of $N_n(\oo)$ lie in $\oo$, where $\psi=1$, we
have $\psi^{-1}(u'^{-1}u)=1$, i.e.\ $\psi^{-1}(u)=\psi^{-1}(u')$. Hence the value
$\psi^{-1}(u)$ is independent of the factorization.

\emph{$(N_n,\psi^{-1})$-equivariance.} For $u_0\in N_n$ and $g=u\,p\in
N_n(P_n\cap K)$, $u_0g=(u_0u)\,p$, so $F_0(u_0g)=\psi^{-1}(u_0u)=\psi^{-1}(u_0)F_0(g)$;
off $N_n(P_n\cap K)$ both sides vanish (the set $N_n(P_n\cap K)$ is left
$N_n$-stable).

\emph{Smoothness and support.} $F_0$ is right invariant under $P_n\cap K$: for
$p_0\in P_n\cap K$ and $g=u\,p$, $g p_0=u\,(p p_0)$ with $p p_0\in P_n\cap K$, so
$F_0(gp_0)=\psi^{-1}(u)=F_0(g)$; off the support both sides vanish. As $P_n\cap K$
is open in $\GL_n$, $F_0$ is locally constant, i.e.\ smooth. Its support
$N_n(P_n\cap K)$ is compact modulo $N_n$ because $(N_n\cap K)\backslash(P_n\cap K)$
is compact. Thus $F_0\in\Cc(N_n\backslash\GL_n,\psi^{-1})$.

The construction of $F_0$ refers only to $N_n$, $P_n\cap K$ and $\psi$, all of which
are fixed data of $\GL_n(F)$; it does not involve $\pi$ or $Q$.
\end{proof}

Let $f_0\in\Cc(N_{n+1}\backslash P_{n+1},\psi^{-1})$ be the element corresponding to
$F_0$ under Lemma~\ref{lem:Cc-iso}, and choose, by Lemma~\ref{lem:miraPi}, a vector
\[
  W\in\Wmod(\Pi,\psi^{-1})\qquad\text{with}\qquad W|_{P_{n+1}}=f_0 .
\]
By Lemma~\ref{lem:F0wd}, $f_0$ and hence $W$ depend only on $\Pi$ (through the choice
realizing $f_0$). \emph{This is the vector $W$ asserted by the Theorem.}

\subsection{Computation of the integral}

\begin{proposition}\label{prop:collapse}
For $W$ as above and any $V\in\Wmod(\pi,\psi)$,
\begin{equation}\label{eq:collapse}
  \Psi(s,W,V)=\psi(-Q)\cdot M(V),\qquad
  M(V):=\int_{(N_n\cap K)\backslash(P_n\cap K)}V(p)\,dp ,
\end{equation}
which is independent of $s$. Consequently $\Psi(s,W,V)$ is finite for every $s$, and
nonzero for all $s\in\mathbb C$ if and only if $M(V)\neq0$ (note $|\psi(-Q)|=1$).
\end{proposition}

\begin{proof}
By the choice of $W$ and \eqref{eq:twist},
\[
  W(\diag(g,1)u_Q)=f_0(\diag(g,1)u_Q)=\psi(-Q\,g_{nn})\,F_0(g).
\]
Since $F_0$ is supported on $N_n(P_n\cap K)$, the integrand of $\Psi(s,W,V)$ vanishes
unless $g\in N_n(P_n\cap K)$. The set $N_n\backslash N_n(P_n\cap K)$ is compact, and
a fundamental domain for $N_n$ acting on the left is $(N_n\cap K)\backslash(P_n\cap K)$;
write $g=u\,p$ with $u\in N_n$ and $p\in P_n\cap K$. On this set:
\begin{itemize}
\item $F_0(u\,p)=\psi^{-1}(u)$ by \eqref{eq:F0};
\item $V(u\,p)=\psi(u)\,V(p)$ by $(N_n,\psi)$-equivariance of $V$;
\item $g_{nn}=p_{nn}=1$: indeed $p\in P_n$ has bottom row $e_n^{\mathsf T}$, so its
$(n,n)$-entry is $1$, and $u\in N_n$ has last row $e_n^{\mathsf T}$, so left
multiplication by $u$ does not change the last row; hence the $(n,n)$-entry of
$g=u\,p$ equals that of $p$, namely $1$. Therefore
$\psi(-Q\,g_{nn})=\psi(-Q)$;
\item $|\det g|=|\det p|=1$, because $p\in P_n\cap K\subseteq K$.
\end{itemize}
Hence on the support the integrand equals
$\psi(-Q)\,\psi^{-1}(u)\,\psi(u)\,V(p)\cdot1=\psi(-Q)\,V(p)$, which depends neither
on $u$ nor on $s$, and is a function of the class of $p$ in
$(N_n\cap K)\backslash(P_n\cap K)$ alone. We now reduce the integral over
$N_n\backslash\GL_n$ to one over this compact quotient. Since the integrand is
supported on $N_n(P_n\cap K)$,
\[
  \Psi(s,W,V)=\int_{N_n\backslash N_n(P_n\cap K)}\psi(-Q)\,V(p)\,d\dot g .
\]
The inclusion $P_n\cap K\hookrightarrow N_n(P_n\cap K)$ induces a homeomorphism
\[
  (N_n\cap K)\backslash(P_n\cap K)\;\xrightarrow{\ \sim\ }\;
  N_n\backslash N_n(P_n\cap K),
\]
because $N_n\cap(P_n\cap K)=N_n\cap K$ and every element of $N_n(P_n\cap K)$ is
$N_n$-equivalent to an element of $P_n\cap K$. Transporting the invariant measure
$d\dot g$ on $N_n\backslash\GL_n$ through this homeomorphism yields a positive Radon
measure $d\mu$ on the compact space $(N_n\cap K)\backslash(P_n\cap K)$, and
\[
  \Psi(s,W,V)=\psi(-Q)\int_{(N_n\cap K)\backslash(P_n\cap K)}V(p)\,d\mu(p).
\]
We normalize the Haar measure on $\GL_n$, hence the quotient measure $d\dot g$, so
that $d\mu$ is the quotient of the chosen Haar measure on $P_n\cap K$ by that on
$N_n\cap K$; this is exactly the measure $dp$ used in the definition of $M(V)$, so
\eqref{eq:collapse} holds. (Any other normalization changes $\Psi$ by a fixed
positive constant and affects none of the conclusions.) The function $V$ is locally
constant on the compact set $(N_n\cap K)\backslash(P_n\cap K)$, so $M(V)$ is a finite
complex number and $\Psi(s,W,V)=\psi(-Q)\,M(V)$ for every $s$. As $|\psi(-Q)|=1\neq0$,
the last assertion follows.
\end{proof}

We have reduced the Theorem to the following statement, which involves only the
mirabolic restriction of $V$:
\begin{equation}\label{eq:goal}
  \boxed{\ \exists\,V\in\Wmod(\pi,\psi)\ \text{with}\
  M(V)=\int_{(N_n\cap K)\backslash(P_n\cap K)}V(p)\,dp\neq0.\ }
\end{equation}

\section{Non-vanishing of $M$ and conclusion}

The non-vanishing \eqref{eq:goal} is now immediate from the standard surjectivity of
the Kirillov model onto the compactly induced part, with \emph{no} prescription of
$V$ outside the mirabolic and \emph{no} use of the conductor of $\pi$.

\begin{lemma}[Kirillov model of $\pi$]\label{lem:miraPiSmall}
Let $\pi$ be generic irreducible admissible on $\GL_n(F)$. The restriction map
$V\mapsto V|_{P_n}$ from $\Wmod(\pi,\psi)$ to functions on $P_n$ is injective, and
its image contains $\Cc(N_n\backslash P_n,\psi)$. In particular, for every
$\phi_0\in\Cc(N_n\backslash P_n,\psi)$ there is $V_0\in\Wmod(\pi,\psi)$ with
$V_0|_{P_n}=\phi_0$.
\end{lemma}

\begin{proof}
This is the same Bernstein--Zelevinsky/Kirillov statement as Lemma~\ref{lem:miraPi},
now for $\GL_n$: the top derivative of $\pi$ realizes
$\operatorname{ind}_{N_n}^{P_n}\psi=\Cc(N_n\backslash P_n,\psi)$ inside
$\Wmod(\pi,\psi)|_{P_n}$, and the restriction to $P_n$ is injective by irreducibility
and genericity of $\pi$. References as in Lemma~\ref{lem:miraPi}.
\end{proof}

\begin{proposition}\label{prop:nonvanish}
There exists $V\in\Wmod(\pi,\psi)$ with $M(V)\neq0$.
\end{proposition}

\begin{proof}
Define $\phi_0$ on $P_n$ by
\[
  \phi_0(u\,p):=\psi(u)\qquad(u\in N_n,\ p\in P_n\cap K),\qquad
  \phi_0\equiv0\ \text{off}\ N_n(P_n\cap K).
\]
This is well defined and lies in $\Cc(N_n\backslash P_n,\psi)$ by exactly the
argument of Lemma~\ref{lem:F0wd} (with $\psi$ in place of $\psi^{-1}$ and $P_n$ in
place of $\GL_n$): if $u\,p=u'\,p'$ then $u'^{-1}u\in N_n\cap K=N_n(\oo)$, on which
$\psi=1$, so the value is well defined; $\phi_0$ is $(N_n,\psi)$-equivariant, right
$(P_n\cap K)$-invariant hence smooth, and supported on $N_n(P_n\cap K)$, which is
compact modulo $N_n$.

By Lemma~\ref{lem:miraPiSmall} choose $V\in\Wmod(\pi,\psi)$ with $V|_{P_n}=\phi_0$.
Since $P_n\cap K\subseteq P_n$, for $p\in P_n\cap K$ we have $V(p)=\phi_0(p)=\psi(1)=1$
(taking $u=I$). Hence
\[
  M(V)=\int_{(N_n\cap K)\backslash(P_n\cap K)}V(p)\,dp
      =\int_{(N_n\cap K)\backslash(P_n\cap K)}1\,dp
      =\operatorname{vol}\bigl((N_n\cap K)\backslash(P_n\cap K)\bigr)>0 .
\]
In particular $M(V)\neq0$.
\end{proof}

\begin{proof}[Proof of the Theorem]
Fix $\Pi$ and let $W\in\Wmod(\Pi,\psi^{-1})$ be the vector constructed in
\S\ref{sec:chosenW}; by Lemma~\ref{lem:F0wd} it depends only on $\Pi$. Let $\pi$ be
an arbitrary generic irreducible admissible representation of $\GL_n(F)$, with
conductor ideal $\qq$ and $Q$ a generator of $\qq^{-1}$. By
Proposition~\ref{prop:nonvanish} there is $V\in\Wmod(\pi,\psi)$ with $M(V)\neq0$. By
Proposition~\ref{prop:collapse},
\[
  \Psi(s,W,V)=\psi(-Q)\,M(V)
\]
is independent of $s$, and $|\psi(-Q)|=1$, $M(V)\neq0$, so it is a fixed
nonzero complex number. Thus $\Psi(s,W,V)$ is finite and nonzero for every
$s\in\mathbb C$. This is precisely the asserted property of $W$.
\end{proof}

\section*{Remarks}

\begin{enumerate}
\item \emph{Why the conductor plays no explicit role.} The element
$u_Q=I_{n+1}+Q\,E_{n,n+1}$ acts on the integrand through the $(n,n)$-entry $g_{nn}$
only, via the factor $\psi(-Q\,g_{nn})$ (Lemma~\ref{lem:Cc-iso}). By supporting the
fixed function $F_0$ on the integral mirabolic $N_n(P_n\cap K)$ --- where every
element has $g_{nn}=1$ --- this factor is frozen to the nonzero constant $\psi(-Q)$.
Thus the single vector $W$, built once from $\Pi$, neither knows nor needs the
conductor exponent $c=v(\qq)$; the dependence of $u_Q$ on $\pi$ is absorbed into an
inessential unimodular constant. This is what makes the choice of $W$ uniform in
$\pi$.

\item \emph{Why no prescription of $V$ on $K$ is needed.} Because the collapsed
integral $M(V)$ only involves the restriction of $V$ to the mirabolic subgroup
$P_n$ (indeed to $P_n\cap K$), the non-vanishing is settled by the single genuinely
available freedom in the Kirillov model: the image of $V\mapsto V|_{P_n}$ contains
all of $\Cc(N_n\backslash P_n,\psi)$ (Lemma~\ref{lem:miraPiSmall}). We do not, and
need not, prescribe $V$ on the part of $K$ lying outside $P_n$; the achievable
restrictions of Whittaker functions to all of $K$ are constrained by the $K$-type
content of $\pi$, but those constraints are irrelevant here.

\item \emph{The integral is constant in $s$.} On the support of the integrand
$|\det g|\equiv1$, so $\Psi(s,W,V)$ is independent of $s$; it is therefore a nonzero
element of $\mathbb C$ with neither zeros nor poles, which is the precise meaning of
``finite and nonzero for all $s\in\mathbb C$.'' (A bare local $L$-factor would not
qualify, as it has poles.)

\item No hypothesis beyond genericity/admissibility of $\pi$ and $\Pi$ is used, and
the construction is uniform in $\pi$ once $W$ is fixed.
\end{enumerate}

\end{document}
